\section*{Bab \ref{ch:b3:petrological-thermodynamics} --- Termodinamika Batuan}

\begin{solution}{exo:b3:petrological-thermodynamics:1}
$\Delta_rH = -3029.3 - 910.7 + 3935.0 = \qty{-5.0}{kJ}$; $\Delta_rS = 133.5 +
41.5 - 207.4 = \qty{-32.4}{J.K^{-1}}$; $\Delta_rV = 60.40 + 22.69 - 100.07 =
\qty{-16.98}{cm^3} = \qty{-1.698}{J.bar^{-1}}$. Kemiringan $32.4/1.698 =
\qty{19.1}{bar/K}$. Pada \qty{873}{K}: $P = 1 + (873.15\times(-32.4) +
5000)/(-1.698) = \qty{13.7}{kbar}$.
\end{solution}

\begin{solution}{exo:b3:petrological-thermodynamics:2}
$c = 2$ dan $p = 3 > c$: aturan fasa dilanggar. Ketiga fasa hanya berada bersama pada
garis univarian \ce{Mg2SiO4 + SiO2 <=> 2 MgSiO3}; batuan yang diambil
secara acak tidak berada di sana, sehingga asosiasi mineral mencatat reaksi yang belum lengkap, misalnya
forsterit yang terselubung tepi enstatit yang memisahkannya dari
kuarsa.
\end{solution}

\begin{solution}{exo:b3:petrological-thermodynamics:3}
$a_{\text{grosular}} = 0.20^3 = 0.008$; $a_{\text{piropa}} = 0.15^3 =
\num{3.4e-3}$; $a_{\text{anortit}} = 0.40$.
\end{solution}

\begin{solution}{exo:b3:petrological-thermodynamics:4}
$T_c = 15\,000/(2\times8.314) = \qty{902}{K}$; $W = 2\times8.314\times1000
= \qty{16.6}{kJ.mol^{-1}}$.
\end{solution}

\begin{solution}{exo:b3:petrological-thermodynamics:5}
$\Delta_rG = -5432.2 + 2\times(-2088.8) - 730.2 + 3\times3478.1 =
\qty{94.3}{kJ}$; $P = 1 + 94\,300/6.61 = \qty{14.3}{kbar}$. Garis eksperimen
memberikan $22.80\times726.85 - 1093 = \qty{15.5}{kbar}$, \qty{1.2}{kbar}
lebih tinggi.
\end{solution}

\begin{solution}{exo:b3:petrological-thermodynamics:6}
$RT = \qty{7483}{J.mol^{-1}}$. Pada $x = 0.5$: $7483\ln0.5 + 20\,000/4 = -5187 +
5000 = \qty{-187}{J.mol^{-1}}$. Pada $x = 0.111$: $7483\times(0.111\ln0.111 +
0.889\ln0.889) + 20\,000\times0.111\times0.889 = -2609 + 1974 =
\qty{-635}{J.mol^{-1}}$, dan nilai yang sama pada 0.889 karena simetri. Jumlah yang sama dari
kedua sayap memiliki komposisi total 0.5 dan energi Gibbs
\qty{-635}{J.mol^{-1}}: pemisahan menurunkannya sebesar \qty{448}{J.mol^{-1}}.
\end{solution}

\begin{solution}{exo:b3:petrological-thermodynamics:7}
$c + 2 = 5$ fasa dan 5 garis. (D) $A + C = 2B$: $0 + 2/3 = 2\times1/3$. (A)
$B + D = 2C$: $1/3 + 1 = 2\times2/3$. (C) $2A + D = 3B$: $0 + 1 = 3\times1/3$.
(B) $A + 2D = 3C$: $0 + 2 = 3\times2/3$.
\end{solution}

\begin{solution}{exo:b3:petrological-thermodynamics:8}
Menurut aturan tuas, fasa pada 0.889 memberikan $(0.3 - 0.111)/(0.889 - 0.111) =
0.243$, \qty{24}{\percent}, dan fasa pada 0.111 \qty{76}{\percent}.
\end{solution}

\begin{solution}{exo:b3:petrological-thermodynamics:9}
$\partial P/\partial X_{\text{an}} = -3RT/(6.21X_{\text{an}}) =
-3\times8.314\times900/(6.21\times0.30) = \qty{-12000}{bar}$ per satuan
$X_{\text{an}}$: kesalahan 0.01 menggeser tekanan sebesar \qty{120}{bar}.
\end{solution}

\begin{solution}{exo:b3:petrological-thermodynamics:10}
$G'_{\text{mix}} = RT\ln[x/(1 - x)] + W(1 - 2x)$ dan $G''_{\text{mix}} =
RT/[x(1 - x)] - 2W$, nol ketika $x(1 - x) = RT/(2W)$. Untuk $x = 0.3$: $T =
2\times20\,000\times0.21/8.314 = \qty{1010}{K}$; \omterm{def:b3:petrological-thermodynamics:solvus}{solvus} tercapai pada
$20\,000\times0.4/(8.314\ln(7/3)) = \qty{1136}{K}$. Di antara keduanya, kristal
bersifat metastabil: pemisahan akan menurunkan energinya, tetapi fluktuasi kecil
menaikkannya, sehingga fasa kedua harus bernukleasi.
\end{solution}

\begin{solution}{exo:b3:petrological-thermodynamics:11}
Karena $G_{\text{mix}}(x) = G_{\text{mix}}(1 - x)$, garis singgung bersama menyentuh
kurva pada dua titik simetris $x$ dan $1 - x$ dan mendatar:
$G'_{\text{mix}}(x) = 0$, yang merupakan persamaannya. Jika diselesaikan, $T(x) = W(1 -
2x)/[R\ln((1 - x)/x)]$. Dengan $x = 1/2 - \varepsilon$, $1 - 2x = 2\varepsilon$
dan $\ln((1/2 + \varepsilon)/(1/2 - \varepsilon)) \approx 4\varepsilon$, sehingga
$T \to W/(2R)$ saat $\varepsilon \to 0$; dan untuk $x < 1/2$, $T(x)$ naik seiring $x$
mendekati 1/2, sehingga limit ini merupakan nilai maksimum.
\end{solution}

\begin{solution}{exo:b3:petrological-thermodynamics:12}
$(2.0 - 0.08\times10)/0.92 = \qty{1.30}{wt\percent}$. Mangan menyukai
garnet: dengan jumlah yang lebih sedikit dalam \omterm{def:b3:petrological-thermodynamics:pseudosection}{komposisi total efektif}, tepi tumbuh dalam
medan yang berbeda dan sepanjang isoplet yang berbeda daripada
yang diperkirakan oleh analisis batuan utuh.
\end{solution}

\begin{solution}{pb:b3:petrological-thermodynamics:1}
\textbf{1.} $\Delta_rH^\circ = -6640.0 + 2\times(-2593.8) - 910.7 + 3\times4234.0
= \qty{-36.3}{kJ}$; $\Delta_rS^\circ = 260.1 + 165.6 + 41.5 - 597.9 =
\qty{-130.7}{J.K^{-1}}$; $\Delta_rV = 125.28 + 88.30 + 22.69 - 302.37 =
\qty{-66.10}{cm^3} = \qty{-6.61}{J.bar^{-1}}$.

\textbf{2.} $-36\,300 + 298.15\times130.7 = \qty{+2.67}{kJ}$: positif, sehingga
anortit stabil di permukaan.

\textbf{3.} $P = 1 + (900\times(-130.7) + 36\,300)/(-6.61) = \qty{12.3}{kbar}$;
kemiringan $130.7/6.61 = \qty{19.8}{bar/K}$.

\textbf{4.} $1 + 81\,500/6.61 = \qty{12.3}{kbar}$ pada \qty{900}{K}; $1 +
119\,800/6.61 = \qty{18.1}{kbar}$ pada \qty{1200}{K}.

\textbf{5.} Pada 626.85 dan \qty{926.85}{\celsius} garis eksperimen memberikan
13.2 dan \qty{20.0}{kbar}: 0.9 dan \qty{1.9}{kbar} di atas tabel.

\begingroup\emergencystretch=3em\relax \textbf{6.} $1914\times6.61 = \qty{12.7}{kJ}$. Ketidakpastian gabungan dari keempat suku ini adalah\linebreak
$\sqrt{3.2^2 + (3\times2.0)^2 + (2\times2.0)^2 + 1.0^2} = \qty{8.0}{kJ}$: selisihnya
1.6 kali itu, mungkin untuk kalorimetri, dan itulah alasan garis ini
diambil dari eksperimen.
\par\endgroup
\textbf{7.} $\Delta_rC_p$ kecil di antara padatan, dan kedua integralnya sebagian
saling meniadakan: tabel hanya menggeser garis sebesar 26 dan \qty{112}{bar}.

\textbf{8.} \ce{SiO2}, \ce{Al2O3}, \ce{FeO}, \ce{MgO}, \ce{MnO}, \ce{CaO},
\ce{Na2O}, \ce{K2O} dan \ce{H2O}: $c = 9$. Fasa: garnet, biotit,
muskovit, plagioklas, kianit, kuarsa dan fluida, $p = 7$.

\textbf{9.} $F = 9 - 7 + 2 = 4$: tekanan, temperatur dan dua
variabel komposisi. $p = 7 \le c = 9$: aturan ini berlaku, dan asosiasi mineral dapat stabil
pada suatu area di bidang tekanan--temperatur.

\textbf{10.} Dua polimorf \ce{Al2SiO5} murni dalam kesetimbangan terletak pada
garis univarian kianit = silimanit: satu hubungan antara $P$ dan $T$, dan
termometer kemudian memberikan tekanannya.

\textbf{11.} Fluida menentukan potensial kimia air: 8 komponen bebas
untuk 6 fasa padat, $F = 8 - 6 + 2 = 4$, dan $p \le c$ tetap
berlaku.

\textbf{12.} $a_{\text{grosular}} = 0.10^3 = 0.001$, $a_{\text{anortit}} =
0.30$; $\ln K = \ln(0.001/0.30^3) = -3.30$.

\textbf{13.} $K_D = (0.12/0.74)/(0.45/0.55) = 0.1622/0.8182 = 0.198$.

\textbf{14.} $T = (52\,635 + 2488)/(20.20 + 24.94\times1.618) = 55\,123/60.57
= \qty{910}{K}$, atau \qty{637}{\celsius}.

\textbf{15.} $P_0 = 22.80\times636.9 - 1093 = \qty{13429}{bar}$; $(3\times8.314
\times910.1/6.21)\ln(1/3) = \qty{-4016}{bar}$; $P = \qty{9.41}{kbar}$.

\textbf{16.} Pada \qty{9.41}{kbar}, \qty{644.2}{\celsius}, lalu
\qty{9.55}{kbar}; pada tekanan itu \qty{644.9}{\celsius} dan
\qty{9.56}{kbar}, tidak berubah setelahnya: \qty{645}{\celsius} dan
\qty{9.56}{kbar}.

\textbf{17.} \qty{641}{\celsius} dan \qty{8.79}{kbar}. Tabel menempatkan
garis komponen akhir sekitar \qty{1}{kbar} terlalu rendah (Bagian I); garis
eksperimen dipertahankan.

\textbf{18.} Dengan $X_{\text{an}} = 0.35$: \qty{642}{\celsius} dan
\qty{8.93}{kbar}, \qty{0.63}{kbar} lebih rendah.

\textbf{19.} Pada \qty{918}{K}: $P = 1 + (918\times12.6 - 7700)/0.571 =
\qty{6.8}{kbar}$, di bawah \qty{9.56}{kbar}: kianit, seperti yang teramati.

\textbf{20.} $z = \num{9.56e8}/(2800\times9.81) = \qty{34.8}{km}$.

\textbf{21.} $\Delta P = (3\times8.314\times918/6.21)\ln(0.07/0.10) =
\qty{-1.3}{kbar}$: inti tumbuh sekitar \qty{4.8}{km} lebih dangkal. Garnet
tumbuh selama penguburan, dan tepinya merekam titik terdalam.

\textbf{22.} Sekis Varden mencapai tekanan puncak sebesar
\textbf{\qty[text-series-to-math=true]{9.6}{kbar}} pada
\qty{645}{\celsius}, dengan penguburan sekitar \qty{35}{km}.
\end{solution}
